\chapter{Einrichtung (einmalig)}
\label{ch:einrichtung}

\section{Python-Umgebung}

\paragraph{Mit dem Setup-Fenster (empfohlen).} Im Repository der Gruppe
\pfad{code\setup_m4.bat} doppelklicken. Das Fenster zeigt jede Komponente mit Status und legt
ohne Administratorrechte \pfad{code\m4_behavior_cloning\.venv} an: Python~3.12, torch (mit
NVIDIA-Karte die CUDA-Version, sonst CPU) und LeRobot. Treiber und Kamera-SDK prüft es nur und
gibt dafür eine Checkliste für den Laborbetreuer aus.

\paragraph{Von Hand.} Die Reihenfolge ist wichtig, sonst überschreibt LeRobot die CUDA-Version
von torch. Ausführlich in \pfad{Dokumentation/Environment-Einrichtung.md}:
\begin{enumerate}
  \item Umgebung mit Python~3.12 anlegen (3.11 geht nicht, LeRobot verlangt 3.12).
  \item torch und torchvision \emph{zuerst} aus dem cu128-Index installieren
    (RTX-50-Karten brauchen sm\_120).
  \item \pfad{python tools/check_gpu.py}: \texttt{sm\_120} muss in der Architekturliste stehen.
  \item \pfad{lerobot[dataset,diffusion]==0.6.1} installieren.
  \item OpenCV auf die volle Variante umstellen (nicht headless).
  \item \pfad{pip install -e .[dev,gui]} für die übrigen Abhängigkeiten.
\end{enumerate}
Die gepinnten Versionen stehen in \pfad{Dokumentation/environment-lock.txt} (PC) und
\pfad{environment-lock-laptop.txt}.

\hinweis{Für das Training ohne Internet die ImageNet-Gewichte von ResNet18 vorher nach
\pfad{~/.cache/torch/hub/checkpoints/resnet18-f37072fd.pth} kopieren.}

\section{Daheng-Kameras (Galaxy SDK)}

\begin{enumerate}
  \item Galaxy SDK des Herstellers mit Administratorrechten installieren (Treiber).
  \item Die Python-Schnittstelle \pfad{gxipy} ist kein pip-Paket. Der Adapter findet sie im
    Installationsordner selbst; abweichend über die Umgebungsvariable
    \pfad{GALAXY_SDK_PYTHON}.
  \item Nach der Installation die Shell neu öffnen. Fehlermeldungen wie
    \pfad{KeyError: 'GALAXY_GENICAM_ROOT'} bedeuten nur, dass die Umgebungsvariablen noch nicht
    angekommen sind.
  \item Prüfen: \pfad{python tools/check_cameras.py --list} zeigt Modell und Seriennummer.
\end{enumerate}

\achtung{Beide Kameras direkt an getrennte USB-3-Anschlüsse des Rechners stecken, nicht an
einen Verteiler. Am Verteiler liefern sie nur etwa 23 statt 60 Bilder pro Sekunde, und die
Episoden werden verworfen.}

\section{Netzwerk}
\label{sec:netzwerk}

\begin{table}[htbp]
  \centering
  \small
  \begin{tabularx}{\textwidth}{@{}l l X@{}}
    \toprule
    Gerät & Adresse & Bemerkung \\
    \midrule
    Steuerung der Anlage & 192.168.2.13, Port 65432 & NeuraPy-Server, Steuerung v5.0.8 \\
    Laptop an der Anlage & 192.168.2.14 & feste IP am LAN-Anschluss, Maske 255.255.255.0 \\
    \acs{vm} \enquote{Lara5\_V5} & 192.168.2.13 & gleiche Adresse wie die Anlage! \\
    Laptop zur \acs{vm} & 192.168.2.11 & VirtualBox Host-Only-Adapter \#3 \\
    \bottomrule
  \end{tabularx}
\end{table}

\begin{itemize}
  \item LAN-Kabel von der Buchse \enquote{Ethernet RJ45} in der Control Box direkt an den Laptop.
  \item IPv4 am Laptop manuell setzen, Gateway und DNS leer lassen. Ist das WLAN selbst ein
    192.168.2.x-Netz, WLAN ausschalten.
  \item Prüfen:
\begin{lstlisting}[style=ps]
ping 192.168.2.13
Test-NetConnection 192.168.2.13 -Port 65432
\end{lstlisting}
\end{itemize}

\achtung{\acs{vm} und Anlage hören auf dieselbe Adresse. Welche gerade antwortet, zeigt
\emph{nur} die Statusanzeige der Oberfläche (rot = reale Anlage, \cref{sec:backend}). Nie
annehmen, mit der \acs{vm} verbunden zu sein. Vor der Arbeit an der Anlage die \acs{vm}
ausschalten.}

\hinweis{In den Unterlagen stehen mehrere Adressen: Die Neura-Anleitung von Modul~2 nennt
192.168.2.20 (dort selbst als \enquote{prüfen} markiert), Neura-Standard ist 192.168.2.13. Mit
Modul~4 lief die Verbindung zur Anlage über 192.168.2.13. Antwortet sie dort nicht, die Adresse
am Touchpad in den Netzwerkeinstellungen nachsehen und
\pfad{Test-NetConnection <Adresse> -Port 65432} prüfen.}

\paragraph{Virtuelle Steuerung.} Die \acs{vm} \enquote{Lara5\_V5} läuft in VirtualBox. Der zweite
Netzwerkadapter der \acs{vm} muss auf \enquote{VirtualBox Host-Only Ethernet Adapter \#3} zeigen;
steht er auf dem Adapter ohne Nummer, startet die \acs{vm} mit
\pfad{VERR_INTNET_FLT_IF_NOT_FOUND}. Der Laptop muss mit Netzteil und im Energiemodus
\enquote{Beste Leistung} laufen, sonst hält er den Takt nicht.
